\documentclass[10pt,a4paper]{amsart}
\usepackage[margin=19mm]{geometry}
\usepackage{amsmath,amssymb,mathtools}
\usepackage[scr=rsfs,cal=euler]{mathalfa}
\usepackage{xfrac}
\usepackage[unicode,hidelinks,bookmarks=false]{hyperref}
\newcommand{\ie}{i.e.\ }
\newcommand{\cf}{cf.\ }
\newcommand{\resp}{respectively\ }
\newcommand{\Subsection}{\subsection}
\providecommand{\nobreakdash}{\nobreak\hbox{-}}
\setlength{\emergencystretch}{2em}
\multlinegap0pt
\usepackage{bm}
\usepackage{mathtools}
\usepackage{xcolor}
\usepackage{tikz}
\usepackage{dsfont}
\usepackage{amssymb}
\usepackage{enumerate}
\usepackage{esint}
\newtheorem{definition}{Definition}
\newtheorem{lemma}{Lemma}
\newtheorem{theorem}{Theorem}
\newcommand{\R}{\ensuremath{\mathbb{R}}}
\DeclareMathOperator{\Tail}{Tail}
\renewcommand{\d}{\:\! \mathrm{d}}
\newcommand{\dx}{\mathrm{d}x}
\renewcommand{\epsilon}{\varepsilon}
\DeclareMathOperator{\loc}{loc}
\renewcommand{\rho}{\varrho}
\newcommand\sfE{{\boldsymbol{\mathsf E}}}
\newcommand\sfI{{\boldsymbol{\mathsf I}}}
\newcommand\sfK{{\boldsymbol{\mathsf K}}}
\newcommand{\spt}{\operatorname{spt}}
\title{Sharp gradient integrability for\\ $(s,p)$-Poisson type equations}
\author{Verena B\"ogelein, Frank Duzaar, Naian Liao, Kristian Moring}
\date{Source: arXiv:2602.08944v1 (CC BY 4.0)}
\begin{document}
\maketitle

\noindent\emph{Source and licence.} Sharp gradient integrability for\\ $(s,p)$-Poisson type equations. Version arXiv:2602.08944v1 (CC BY 4.0), distributed by arXiv under the Creative Commons Attribution 4.0 International licence, http://creativecommons.org/licenses/by/4.0/, as recorded in the arXiv licence metadata for this exact version and shown on the abstract page https://arxiv.org/abs/2602.08944v1. Full licence notice: \texttt{licenses/arxiv-2602.08944-CC-BY-4.0.txt}.\par\smallskip

\noindent\textit{Editorial scope.} Counted unit: the article's theorem thm:higher-diff (source0.tex lines 2006--2035). The article's complete original proof of it is delivered with it (source0.tex lines 3249--3435, one \texttt{proof} environment of 187 lines, which the article places away from the statement). No mathematics is written, repaired or re-derived: the statement and the proof are the article's own lines, in the article's own notation, and no retained line is substituted or re-flowed.  Retained uncounted as the source-local context the proof needs: lines 253--255; lines 261--268; lines 483--500; lines 514--521; lines 524--531; lines 1429--1437; lines 1538--1540; lines 1562--1573; lines 1578--1583; lines 1772--1779; lines 1875--1908; lines 2107--2112; lines 2114--2121; lines 2126--2140; lines 2183--2187; lines 2395--2400; lines 2739--2746; lines 2751--2789; lines 2974--3012.  Nothing was omitted from any retained span, so the package carries no omission record.  No retained line contains a citation, so the wrapper carries no bibliography and no bibliography entry.  No retained line contains a figure environment, an image-inclusion command or drawing code, so no image is delivered and the package declares no asset dependency.  Editorial formatting only, in addition: an AMS \texttt{amsart} preamble that restates the article's own declarations and macros used by the retained lines, namely the environments \texttt{definition}, \texttt{lemma}, \texttt{theorem} on separate counters, together with the standard packages those lines need.  The retained environments print in this document as Definition 1, Lemma 1, Theorem 1 in order of appearance, whereas the article prints the same items with the numbering of its own full document.  The complete native source of the article as distributed in the arXiv e-print for this exact version is bundled unchanged as \texttt{sources/194284/source0.tex}.
\begin{equation} \label{eq:frac-p-lap}
\left( - a \Delta_p\right)^s u = f \quad \text{ in } \Omega,
\end{equation}

\begin{equation}\label{eq:a-condition}
\left\{
    \begin{array}{cc}
       C_o\le a(x,y)\le C_1 \quad  \mbox{a.e. $(x,y)\in\R^n\times\R^n$} \\[5pt]
       \displaystyle\sup_{x,y\in B_R(x_o)}|a(x,y)-a(x_o,x_o)|\le C_1 R^{\chi}\quad \mbox{for any $R\in(0,R_o]$}%x\mapsto a(x,x)\in C^{\chi}_{\loc}(\Omega)   
    \end{array}
    \right.
\end{equation}

\begin{definition}\label{def:weak-sol}
Let $\Omega\subset\mathbb R^n$ be a bounded open set, $p\in(1,\infty)$, $s\in(0,1)$, and $f\in L^{r}_{\rm loc}(\Omega)$, with
$$
    r\ge (p_s^\ast)'\quad\mbox{if $sp\not=n$,}
    \qquad\mbox{or}\qquad 
    r>1 \quad\mbox{if $sp=n$.}
$$
A function $u\in W^{s,p}_{\rm loc}(\Omega)$
is called a  {\bf local weak solution} to the fractional $p$-Poisson equation \eqref{eq:frac-p-lap} in $\Omega$ if and only if
\begin{equation}\label{Lebesgue-weight}
\int_{\mathbb R^n}\frac{|u(x)|^{p-1}}{(1+|x|)^{n+sp}}\, \mathrm{d}x <\infty,
\end{equation}
and
\begin{equation*}
        \iint_{\mathbb R^n\times\mathbb R^n}a(x,y)\frac{|u(x)-u(y)|^{p-2}(u(x)-u(y))(\varphi(x)-\varphi(y))}{|x-y|^{n+sp}}\,\mathrm{d}x \mathrm{d}y =2\int_{\Omega} f \varphi \, \d x
\end{equation*}
for every $\varphi\in W^{s,p}(\R^n)$ with $\spt(\varphi)\Subset\Omega$.
\end{definition}

\begin{equation}\label{Dirichlet} 
\left\{
    \begin{array}{cl}
      (- a\Delta_p)^s u  = 0 & \mbox{in $\Omega$,} \\[6pt]
      u = g & \mbox{a.e.~in $\R^n \setminus \Omega$.}
    \end{array}
\right.
\end{equation}

\begin{definition}\label{def:weak-sol-D}
Let $\Omega\subset\mathbb R^n$ be a bounded open set, $p\in(1,\infty)$ and $s\in(0,1)$. A function $u\in W^{s,p}(\Omega)\cap L^{p-1}_{sp}(\R^n)$ is called a {\bf weak solution} to the Dirichlet problem \eqref{Dirichlet} in $\Omega$ if and only if $u=g$ a.e. in $\R^n\setminus\Omega$
and
\begin{equation*}
        \iint_{\mathbb R^n\times\mathbb R^n}a(x,y)\frac{|u(x)-u(y)|^{p-2}(u(x)-u(y))(\varphi(x)-\varphi(y))}{|x-y|^{n+sp}}\,\mathrm{d}x \mathrm{d}y =0
\end{equation*}
for every $\varphi\in W^{s,p}(\R^n)$ satisfying $\varphi\equiv0$ a.e. in $\R^n\setminus\Omega$.
\end{definition}

\begin{equation}\label{def:theta}
    \theta :=
    \left\{
    \begin{array}{cl}
        p+\frac12 sp, & \text{if } p\in(1,2),\\[3pt]
        sp+2, & \text{if } p\ge2.
    \end{array}
    \right.
\end{equation}

\begin{equation} \label{eq:frac-p}
        \left( - \Delta_p\right)^s u = f \quad \text{ in } \Omega,
\end{equation}

\begin{align}\label{def:alpha}
    \alpha 
    &:=
    \max\bigg\{\frac{n}{(p-1)\,[\,n+p(sp'-1)\,]}\,,\, \frac{1}{p}\bigg\} \nonumber \\[6pt]
    &=\left\{
    \begin{array}{cl}
    \displaystyle
        \frac{n}{(p-1)[n+p(sp'-1)]}, & \mbox{if $\displaystyle s\in(\tfrac{p-1}{p},s_o]$,}\\[12pt]
        \displaystyle\frac{1}{p}, &  \mbox{if $\displaystyle s\in(s_o,1)$.}
    \end{array}
    \right.
\end{align}

\begin{align}\label{def:eps}
    \varepsilon
    &:=
    sp' - 1 - \frac{n}{p}\bigg[\frac{1}{\tilde\alpha (p-1)}-1\bigg]
    \in [0,\, sp'-1].
\end{align}

\begin{equation}\label{CD-v}
\left\{
    \begin{array}{cl}
      (-\Delta_p)^s v = 0 & \text{in } B_{8R}(x_o), \\[6pt]
      v = u & \text{a.e.\ in } \mathbb{R}^n\setminus B_{8R}(x_o).
    \end{array}
\right.
\end{equation}

\begin{lemma}\label{lem:vi}
Let $p>1$ and $s\in(0,1)$ with $sp'>1$, and let $\tilde\alpha\in[\alpha,\frac{1}{p-1}]$ be such that $\epsilon>0$, where $\alpha$ and $\epsilon$ are defined in~\eqref{def:alpha} and~\eqref{def:eps}. 
Then there exists a constant $c=c(n,p,s,\tilde\alpha)$ such that the following holds.
Whenever
\(u \in W^{s,p}(B_\rho(y_o))\cap L^{p-1}_{sp}(\mathbb{R}^n)\)
is a weak solution in \(B_\rho(y_o)\) satisfying 
\(u \in W^{\sigma,p}(B_\rho(y_o))\) for some \(\sigma\in[s,1)\), 
\(x_o\in B_{\frac12 \rho}(y_o)\),  \(R\in(0,\frac1{16}\rho]\),
and
\(v \in W^{s,p}(B_{8R}(x_o))\cap L^{p-1}_{sp}(\mathbb{R}^n)\)
is the unique weak solution to \eqref{CD-v},
then for any \(m_o\in\mathbb{N}_0\) such that
\(\tfrac{1}{32}\rho<2^{m_o}R\le\tfrac{1}{16}\rho\) and any
step size \(h\in\mathbb{R}^n\) with \(0<|h|\le \tfrac17 R\), there holds
\begin{align*}
    \int_{B_R(x_o)} \big|\boldsymbol{\tau}_h^2 v\big|^p \,\dx
    &\le
    c \Big(\frac{|h|}{R}\Big)^{\theta}
    \Bigg[
        R^{\sigma p}
        \Bigg(
            \sum_{k=0}^{m_o}
            2^{k(\sigma-sp' - \frac{n}{p})}
            [u]_{W^{\sigma,p}(B_{2^k 8R}(x_o))}
        \Bigg)^p \\
    &\qquad\qquad\qquad
        + \Big(\frac{R}{\rho}\Big)^{spp'} R^n\sfE\big(u;B_\rho(y_o)\big)^p \\
    &\qquad\qquad\qquad + 
        R^{(1+\epsilon)p}
        \|f\|_{L^{\tilde\alpha p}(B_{8R}(x_o))}^{p'}
    \Bigg],
\end{align*}
where $\theta$ is given in~\eqref{def:theta}.
\end{lemma}

\begin{align}\label{Eq:tau-h-B-i}
    \int_{B_i} |\boldsymbol\tau_h^2 u|^p \,\d x
    &\le
    2^{p-1}\int_{B_i} |\boldsymbol\tau_h^2 (u - v_i)|^p \,\d x
    + 2^{p-1}\int_{B_i} |\boldsymbol\tau_h^2 v_i|^p \,\d x 
\end{align}

\begin{equation}\label{def:Ii}
    \sfI_i
    :=
    \Big(\frac{|h|}{\rho}\Big)^{sp+(\sigma-s+\bar\epsilon)\beta p} 
    \Big[ \rho^{\sigma p} [u]_{W^{\sigma,p}(16B_i)}^p +
    \rho^{(1+\epsilon)p}
    \|f\|_{L^{\tilde\alpha p}(8B_i)}^{p'}\Big]. 
\end{equation}

\begin{align}\label{change-1}
    \int_{B_i} \big|\boldsymbol\tau_h^2 (u - v_i)\big|^p \,\d x
    &\le
    2^{p-1}\int_{\mathbb{R}^n} \big|\boldsymbol\tau_h (u - v_i)\big|^p \,\d x \nonumber\\
    &\le
    c\, |h|^{sp} [u-v_i]_{W^{s,p}(\mathbb{R}^n)}^p \nonumber\\
    &\le
    c\, |h|^{sp} \big(|h|^\beta\rho^{1-\beta}\big)^{(1-s+\epsilon)p}
    \|f\|_{L^{\tilde\alpha p}(8B_i)}^{p'} \nonumber\\
    &\le 
    c\, \Big(\frac{|h|}{\rho}\Big)^{sp+(\sigma-s+\bar\epsilon)\beta p} \rho^{(1+\epsilon)p}
    \|f\|_{L^{\tilde\alpha p}(8B_i)}^{p'} \nonumber\\
    &\le  
    c\,\sfI_i ,
\end{align}

\begin{align}\label{change-1<}
    \int_{B_i} \big|\boldsymbol\tau_h^2 (u - v_i)\big|^p \,\d x 
    &\le 
    c\,\sfI_i.
\end{align}  

\begin{align}\label{sstep1}
    \int_{B_{\frac12 \rho}(y_o)} \big|\boldsymbol\tau_h^2 u\big|^p \,\dx
    \le
    c\, \Big(\frac{|h|}{\rho}\Big)^{\frac{\sigma(\theta-sp)+\bar\epsilon\theta}{\theta-sp+\bar\epsilon p}\, p}
    \sfK^p,
\end{align}

\begin{equation}\label{CD-homo-a} 
\left\{
    \begin{array}{cl}
      (- a_{x_o,R}\Delta_p)^s v  = 0, & \mbox{in $B_R(x_o)$,} \\[6pt]
      v = u, & \mbox{a.e.~in $\R^n \setminus B_R(x_o)$.}
    \end{array}
\right.
\end{equation}

\begin{lemma}[Comparison estimate with variable coefficients]\label{lem:compar-var-coeff}
Let $p>1$ and $s\in(0,1)$ with $sp'>1$, and let $\tilde\alpha\in[\alpha,\frac{1}{p-1}]$ be such that $\epsilon>0$, where $\alpha$ and $\epsilon$ are defined in~\eqref{def:alpha} and~\eqref{def:eps}.
Then there exists a constant 
$c=c(n,p,s,C_o,C_1,\tilde\alpha)>0$ with the following property. 
Whenever $u$ is a local weak solution to~\eqref{eq:frac-p-lap} in the sense of Definition~\ref{def:weak-sol} with coefficients $a$ satisfying~\eqref{eq:a-condition}, $f\in L^{\tilde\alpha p}(\Omega)$, $B_{2R}\equiv B_{2R}(x_o)\Subset\Omega$ with $R \leq R_o$ where $R_o$ is from~\eqref{eq:a-condition}, and $v\in W^{s,p}(B_R(x_o)) \cap L^{p-1}_{sp}(\R^n)$ denotes the unique weak solution to the homogeneous Dirichlet problem~\eqref{CD-homo-a} with frozen coefficients $a_{x_o,R}$ in the sense of Definition~\ref{def:weak-sol-D}, the function
\[
w := u-v
\]
satisfies the following estimates:


\medskip
\noindent
\emph{(i) Superquadratic case.} 
If $p\in[2,\infty)$, then
\begin{align*}
    R^{-sp}\|w\|_{L^p(B_R)}^p + [w]_{W^{s,p}(\R^n)}^p
    &\leq 
    c\,R^{\frac{\chi p}{p-1}}[u]^p_{W^{s,p}(B_R)} +
    c\, R^{(1-s+\epsilon)p} \|f\|_{L^{\tilde\alpha p}(B_R)}^{p'}.
\end{align*}

\medskip
\noindent
\emph{(ii) Subquadratic case.} 
If $p\in(1,2]$,  then for every $\delta\in(0,1]$ we have
\begin{align*}
    R^{-sp}\|w\|_{L^p(B_R)}^p &+ [w]^p_{W^{s,p}(B_{2R})}\\
    &\le 
    \bigg[\delta 
     +
     \frac{R^{\frac{\chi p}{p-1}}}{\delta^{\frac{2-p}{p-1}p'}} \bigg]
     [u]_{W^{s,p}(B_{2R})}^p 
     + 
     \frac{c}{\delta^{\frac{2-p}{p-1}p'}} 
     R^{(1-s+\epsilon)p}
    \|f\|_{L^{\tilde\alpha p}(B_R)}^{p'} .
\end{align*}
\end{lemma}

\begin{lemma}[Higher differentiability for frozen-coefficient solutions]\label{lem:vi-a}
Let $p>1$, $s\in(0,1)$ with $sp'>1$, and let $B_R\equiv B_R(x_o)\subset\R^n$. Whenever $v\in W^{s,p}_{\rm loc}(B_R)\cap L^{p-1}_{sp}(\R^n)$ is a local weak solution to 
\[
(- a_{x_o,R}\Delta_p)^s v = 0 \quad \text{in } B_R,
\]
then 
\begin{equation*}
    v \in W^{1+\gamma,p}_{\rm loc}(B_R) 
    \quad \text{for every } \gamma \in (0,\gamma_o),
    \qquad 
    \gamma_o := \frac{\theta - p}{\theta - sp + \bar\epsilon p}\,\bar\epsilon,
\end{equation*}
where $\bar\epsilon := \min\{1,p-1\}\,(sp'-1)$ and $\theta$ is defined in~\eqref{def:theta}. 
Moreover, there exists a constant 
$c=c(n,p,s,C_o,C_1,\gamma)>0$ such that for every such $v$ one has
\begin{align*}
    \|\nabla v\|_{L^p(B_{\frac14R})}
    &+   R^\gamma[\nabla v]_{W^{\gamma,p}(B_{\frac14R})} \\
    &\leq
    \frac{c}{R} \Big[ R^{s}[v]_{W^{s,p}(B_{\frac12R})} +
    R^\frac{n}{p} \Tail\big(v-(v)_{\frac12R};B_{\frac12 R}\big)\Big],
\end{align*}
and, for every $h\in\R^n\setminus\{0\}$ with $|h|\le \frac1{32}R$,
\begin{align*}
    \int_{B_{\frac14 R}} |\boldsymbol{\tau}_h^2 v|^p \,\dx
    &\le 
    c \Big(\frac{|h|}{R}\Big)^{(1+\gamma)p} 
    \Big[ R^{sp} [v]_{W^{s,p}(B_{\frac12 R})}^p + R^n \Tail(v-(v)_{\frac12 R}; B_{\frac12 R})^p \Big].
\end{align*}
Moreover, for every $q\in(1,\infty)$, we have $v \in W^{1,pq}_{\rm loc}(B_R)$, and there exists a constant 
$c=c(n,p,s,C_o,C_1,q)>0$ such that
\begin{align*}
    \|\nabla v\|_{L^{pq}(B_{\frac14 R})} 
    &\le 
    c \, R^{-\frac{n}{p}\left(1-\frac{1}{q}\right)-1} 
    \Big[ R^s [v]_{W^{s,p}(B_R)} + 
    R^{\frac{n}{p}} \Tail\big(v-(v)_{\frac12 R}; B_{\frac12 R}\big) \Big].
\end{align*}
\end{lemma}

\begin{theorem}[Higher differentiability]\label{thm:higher-diff}
Let $p>1$ and $s\in(0,1)$ with $sp'>1$, and let $\tilde\alpha\in[\alpha,\frac{1}{p-1}]$ be such that $\epsilon>0$, where $\alpha$ and $\epsilon$ are defined in~\eqref{def:alpha} and~\eqref{def:eps}. Assume in addition that
$f \in L^{\tilde{\alpha}p}_{\mathrm{loc}}(\Omega)$. Then any local weak solution $u$ to~\eqref{eq:frac-p}
belongs to
\begin{equation*}%\label{def:gamma0}
u \in W^{1+\gamma,p}_{\mathrm{loc}}(\Omega)
\quad \text{for every } \gamma \in (0,\gamma_o),
\qquad
\gamma_o := \frac{\theta-p}{\theta-sp+\bar\epsilon p}\,\bar\epsilon,
\end{equation*}
where $\bar\epsilon:=\epsilon \min\{1,p-1\}$ and $\theta$ is given by~\eqref{def:theta}.
Moreover, for every $\gamma\in(0,\gamma_o)$ there exists a constant
$c = c(n,p,s,\tilde{\alpha},\gamma)>0$ such that for every ball
$B_R\Subset\Omega$ the quantitative estimate
\begin{align*}
    &\|\nabla u\|_{L^p(B_{\frac12 R})}
    + R^\gamma [\nabla u]_{W^{\gamma,p}(B_{\frac12 R})}\\
    &\qquad\qquad\le
    \frac{c}{R}
    \bigg[
        R^s [u]_{W^{s,p}(B_R)}
        + R^{\frac{n}{p}}
        \Tail\big(u-(u)_{R}; B_R\big)  
        + R^{1+\epsilon}
        \|f\|_{L^{\tilde{\alpha}p}(B_R)}^{\frac{1}{p-1}}
    \bigg]
\end{align*}
holds. Finally, in the limit $\epsilon\downarrow0$ one has
$\gamma\downarrow 0$ and $c\uparrow\infty$.
\end{theorem}

\begin{proof}
The proof follows the same strategy as the one for the fractional Poisson equation
with constant coefficients $a\equiv1$, cf.\ Theorem~\ref{thm:higher-diff}.
We therefore only indicate the modifications that are required in the present setting.
Throughout the proof, Steps~1--4 refer to the corresponding steps in the proof of
Theorem~\ref{thm:higher-diff}.


\smallskip
\noindent
\textit{Step~1: Estimate for second-order difference quotients.}
We proceed as in the proof of Theorem~\ref{thm:higher-diff} up to
inequality~\eqref{Eq:tau-h-B-i}. 
It remains to estimate the first term on the right-hand side of
\eqref{Eq:tau-h-B-i}, namely
\[
     \int_{B_i} \big|\boldsymbol\tau_h^2 (u - v_i)\big|^p \,\d x,
\]
in terms of the quantity
\[
    \sfI_i
    :=
    \Big(\frac{|h|}{\rho}\Big)^{sp+(\sigma-s+\bar\epsilon)\beta p} 
    \Big[
        \rho^{\sigma p} [u]_{W^{\sigma,p}(16B_i)}^p
        +
        \rho^{(1+\epsilon)p}
        \|f\|_{L^{\tilde\alpha p}(8B_i)}^{p'}
    \Big],
\]
where
\begin{equation}\label{def:bareps}
    \bar\epsilon
    :=
    \min\{1,(p-1)^2\}
    \min\bigg\{\epsilon,\frac{\chi}{p-1}\bigg\}.
\end{equation}
Observe that the quantity $\sfI_i$ coincides with the one defined
in~\eqref{def:Ii}, up to the different definition of $\bar\epsilon$
in Theorem~\ref{thm:higher-diff} and in~\eqref{def:bareps}.
For this estimate, we proceed as in~\eqref{change-1}, applying
the comparison estimate for variable coefficients from
Lemma~\ref{lem:compar-var-coeff}.
In the case $p\ge2$ one has
\begin{align*}
    \int_{B_i} &
    \big|\boldsymbol\tau_h^2 (u - v_i)\big|^p \,\d x\\
    &\le
    c |h|^{sp} [u-v_i]_{W^{s,p}(\mathbb{R}^n)} \\
    &\le
    c |h|^{sp} \Big[
    \big(|h|^\beta\rho^{1-\beta}\big)^{\frac{\chi p}{p-1}}[u]^p_{W^{s,p}(8B_i)} 
    +
    \big(|h|^\beta\rho^{1-\beta}\big)^{(1-s+\epsilon)p}
    \|f\|_{L^{\tilde\alpha p}(8B_i)}^{p'} \Big] \\
    &\le
    c |h|^{sp} \Big[
    \big(|h|^\beta\rho^{1-\beta}\big)^{\frac{\chi p}{p-1}+(\sigma-s)p}[u]^p_{W^{\sigma,p}(8B_i)}  +
    \big(|h|^\beta\rho^{1-\beta}\big)^{(1-s+\epsilon)p}
    \|f\|_{L^{\tilde\alpha p}(8B_i)}^{p'} \Big] \\
    &=
    c \Big(\frac{|h|}{\rho}\Big)^{sp+(\sigma-s)\beta p} 
    \bigg[ 
    \Big(\frac{|h|}{\rho}\Big)^{\frac{\chi}{p-1}\beta p} \rho^{(\sigma+\frac{\chi}{p-1})p} [u]^p_{W^{\sigma,p}(8B_i)} \\
    &\qquad\qquad\qquad\qquad\quad +
    \Big(\frac{|h|}{\rho}\Big)^{(1-\sigma+\epsilon)\beta p} \rho^{(1+\epsilon)p}
    \|f\|_{L^{\tilde\alpha p}(8B_i)}^{p'} \bigg] \\
    &\le
    c \Big(\frac{|h|}{\rho}\Big)^{sp+(\sigma-s)\beta p} 
    \bigg[ 
    \Big(\frac{|h|}{\rho}\Big)^{\frac{\chi}{p-1}\beta p} \rho^{\sigma p} [u]^p_{W^{\sigma,p}(8B_i)} 
    +
    \Big(\frac{|h|}{\rho}\Big)^{\epsilon\beta p} \rho^{(1+\epsilon)p}
    \|f\|_{L^{\tilde\alpha p}(8B_i)}^{p'} \bigg]\\
    &\le 
    c\, \sfI_i,
\end{align*}
where the last inequality follows from the definition of $\bar\epsilon$ in~\eqref{def:bareps} and the assumption $\rho \le R \le 1$. This provides the analogue of~\eqref{change-1} for the present setting.
In the case $p\in(1,2)$, one obtains for any $\delta \in (0,1]$ that
\begin{align*}
    \int_{B_i} \big|\boldsymbol\tau_h^2 (u - v_i)\big|^p \,\d x
    &\le
    c |h|^{sp} [u-v_i]_{W^{s,p}(16B_i)} \\
    &\le
    c |h|^{sp} \Bigg[
    \bigg[\delta +
    \frac{\big(|h|^\beta\rho^{1-\beta}\big)^{\frac{\chi p}{p-1}}}{\delta^{\frac{2-p}{p-1}p'}}\bigg][u]^p_{W^{s,p}(16B_i)} \\
    &\qquad\qquad\quad +
    \frac{\big(|h|^\beta\rho^{1-\beta}\big)^{(1-s+\epsilon)p}}{\delta^{\frac{2-p}{p-1}p'}} 
    \|f\|_{L^{\tilde\alpha p}(8B_i)}^{p'}\Bigg] \nonumber\\
    &\le
    c |h|^{sp} \Bigg[
    \bigg[\delta +
    \frac{\big(|h|^\beta\rho^{1-\beta}\big)^{\frac{\chi p}{p-1}}}{\delta^{\frac{2-p}{p-1}p'}}\bigg]\big(|h|^\beta\rho^{1-\beta}\big)^{(\sigma-s)p} [u]_{W^{\sigma,p}(16B_i)}^p \\
    &\qquad\qquad\quad + 
    \frac{\big(|h|^\beta\rho^{1-\beta}\big)^{(1-s+\epsilon)p}}{\delta^{\frac{2-p}{p-1}p'}}
    \|f\|_{L^{\tilde\alpha p}(8B_i)}^{p'} \Bigg] \\
    &=
    c \Big(\frac{|h|}{\rho}\Big)^{sp+(\sigma-s)\beta p} \Bigg[
    \bigg[\delta +
    \Big(\frac{|h|}{\rho}\Big)^{\frac{\chi}{p-1}\beta p} \frac{\rho^{\frac{\chi p}{p-1}}}{\delta^{\frac{2-p}{p-1}p'}}\bigg]
    \rho^{\sigma p}[u]_{W^{\sigma,p}(16B_i)}^p \\
    &\qquad\qquad\quad + 
    \Big(\frac{|h|}{\rho}\Big)^{(1-\sigma+\epsilon)\beta p} \frac{\rho^{(1+\epsilon)p}}{\delta^{\frac{2-p}{p-1}p'}}
    \|f\|_{L^{\tilde\alpha p}(8B_i)}^{p'} \Bigg].
\end{align*}
We now set 
\[
\delta=\Big(\frac{|h|}{\rho}\Big)^{\bar\epsilon\beta p}.
\]
Then
\begin{align*}
    \delta +
    \Big(\frac{|h|}{\rho}\Big)^{\frac{\chi}{p-1}\beta p} \frac{\rho^{\frac{\chi p}{p-1}}}{\delta^{\frac{2-p}{p-1}p'}}
    =
    \Big(\frac{|h|}{\rho}\Big)^{\bar\epsilon\beta p} +
    \rho^{\frac{\chi p}{p-1}} \Big(\frac{|h|}{\rho}\Big)^{[\frac{\chi}{p-1}-\frac{2-p}{p-1}p'\bar\epsilon]\beta p}
    \le 
    2\Big(\frac{|h|}{\rho}\Big)^{\bar\epsilon\beta p},
\end{align*}
where we used that \(\rho \le R \le 1\) and that
$$
    \frac{\chi}{p-1}-\frac{2-p}{p-1}p'\bar\epsilon
    \ge 
    \frac{\bar\epsilon}{(p-1)^2}-\frac{2-p}{p-1}p'\bar\epsilon
    =
    \frac{\bar\epsilon}{(p-1)^2} \big[1-(2-p)p\big]
    =
    \bar\epsilon
$$
by the definition of \(\bar\epsilon\) in~\eqref{def:bareps}. Moreover, we have 
$$
    \delta^{-\frac{2-p}{p-1}p'}\Big(\frac{|h|}{\rho}\Big)^{(1-\sigma+\epsilon)\beta p}
    \le  
    \Big(\frac{|h|}{\rho}\Big)^{-\frac{2-p}{p-1}p'\bar\epsilon\beta p} \Big(\frac{|h|}{\rho}\Big)^{\epsilon\beta p}
    \le 
    \Big(\frac{|h|}{\rho}\Big)^{\bar\epsilon\beta p},
$$
since \(\sigma<1\) and 
$$
    \epsilon-\frac{2-p}{p-1}p'\bar\epsilon
    \ge 
    \frac{\bar\epsilon}{(p-1)^2}-\frac{2-p}{p-1}p'\bar\epsilon
    =
    \frac{\bar\epsilon}{(p-1)^2} \big[1-(2-p)p\big]
    =
    \bar\epsilon
$$
by the definition of \(\bar\epsilon\) in~\eqref{def:bareps}.
In conclusion, we obtain
\begin{align*}
    \int_{B_i} \big|\boldsymbol\tau_h^2 (u - v_i)\big|^p \,\d x
    &\le
    c \Big(\frac{|h|}{\rho}\Big)^{sp+(\sigma-s+\bar\epsilon)\beta p} 
    \Big[ \rho^{\sigma p}[u]_{W^{\sigma,p}(16B_i)}^p +
    \rho^{(1+\epsilon)p}
    \|f\|_{L^{\tilde\alpha p}(8B_i)}^{p'} \Big] \\
    &=
    c\, \sfI_i .
\end{align*}
Note that this inequality is precisely the analogue of~\eqref{change-1<}.

The estimates for the remaining terms remain unchanged, apart from the choice of $\theta$, 
since Lemma~\ref{lem:vi-a} exactly parallels Lemma~\ref{lem:vi}. 
Hence, we arrive at the analogue of~\eqref{sstep1}, namely 
\begin{align*}
    \int_{B_{\frac12 \rho}(y_o)} \big|\boldsymbol\tau_h^2 u\big|^p \,\dx
    \le
    c \Big(\frac{|h|}{\rho}\Big)^{\frac{\sigma(\theta-sp)+\bar\epsilon\theta}{\theta-sp+\bar\epsilon p}\, p}
    \sfK^p,
\end{align*}
where $c=c(n,p,s,C_o,C_1,\chi,\tilde\alpha,\sigma)$, $\theta=\theta(s,p)$, and
\begin{equation*}
    \sfK^p
    :=
    \rho^{\sigma p}[u]_{W^{\sigma,p}(B_{\rho}(y_o))}^p + 
    \rho^{n}\Tail\big(u-(u)_{y_o,\rho};B_\rho(y_o)\big)^p + \rho^{(1+\epsilon) p}\|f\|_{L^{\tilde\alpha p}(B_\rho(y_o))}^{p'}.
\end{equation*}

\smallskip

\noindent
\textit{Steps 2--4.} 
From this point onward, the proof proceeds exactly as in Theorem~\ref{thm:higher-diff}, 
with the only modification being the definition of $\bar\epsilon$ in~\eqref{def:bareps}. 
Since no further adjustments are needed, we omit the details and refer the reader to Steps 2--4 in the proof of Theorem~\ref{thm:higher-diff}. 
\end{proof}

\end{document}
